\chapter*{Further Reading}
\addcontentsline{toc}{chapter}{Further Reading}
\markboth{Further Reading}{Further Reading}

The References section lists every work cited in the text.
This section points beyond them --- to books and papers worth
reading if the ideas here opened up questions you want to
follow further.

\medskip

\textbf{On data, information, and knowledge.}
Luciano Floridi's \textit{Information: A Very Short Introduction}
(Oxford, 2010) is the most accessible rigorous treatment of what
information actually is.
Russell Ackoff's short 1989 paper ``From Data to Wisdom'' (in the
\textit{Journal of Applied Systems Analysis}) introduced the DIKW
pyramid and remains the standard starting point.

\medskip

\textbf{On representation and semiotics.}
The Stanford Encyclopedia of Philosophy entry on ``Mental
Representation'' (\url{https://plato.stanford.edu/entries/mental-representation/})
gives a careful philosophical treatment.
For semiotics more broadly, Umberto Eco's \textit{A Theory of
Semiotics} (Indiana University Press, 1976) is the classic text.

\medskip

\textbf{On the physical substrate.}
The remaining volumes of \textit{Arliz: Zero to Bit} take up exactly
the question left open at the end of Chapter~3: how an abstract
representation becomes a physical signal. They are the recommended
next step, and are available at the project repository
(\url{https://github.com/papyrxis/arliz}).
